\chapter{Bilirubin (Blood)}

\section*{What Is This Test?}
The bilirubin blood test measures the levels of total bilirubin, direct (conjugated) bilirubin, and indirect (unconjugated) bilirubin in the serum. Bilirubin is the yellow-orange breakdown product of heme, derived primarily from senescent red blood cells (approximately 80\%) and from ineffective erythropoiesis and turnover of other hemoproteins such as myoglobin and cytochromes (approximately 20\%). Heme is converted to biliverdin by heme oxygenase, then to unconjugated bilirubin by biliverdin reductase. Unconjugated bilirubin is water-insoluble and must be transported in plasma tightly bound to albumin. It is taken up by hepatocytes via organic anion transport proteins (OATP1B1 and OATP1B3) and conjugated with glucuronic acid by uridine diphosphate glucuronosyltransferase 1A1 (UGT1A1) to form water-soluble conjugated bilirubin (direct bilirubin). Conjugated bilirubin is excreted into bile canaliculi via the multidrug resistance-associated protein 2 (MRP2/ABCC2) and stored in the gallbladder. In the intestine, conjugated bilirubin is deconjugated by gut bacteria and reduced to urobilinogen, which is either excreted in stool (as stercobilin, giving stool its brown color), reabsorbed into the enterohepatic circulation, or excreted in urine (as urobilin). The differential measurement of total and direct bilirubin allows indirect bilirubin to be calculated (indirect = total -- direct) and is essential for classifying the etiology of jaundice as pre-hepatic (unconjugated hyperbilirubinemia), hepatic (mixed or predominantly conjugated), or post-hepatic/obstructive (conjugated hyperbilirubinemia).

\section*{Alternative Names}
Total Bilirubin; Direct Bilirubin; Conjugated Bilirubin; Indirect Bilirubin; Unconjugated Bilirubin; Serum Bilirubin; T Bili; D Bili; Neonatal Bilirubin; Icterus Index; Fractionated Bilirubin

\section*{Principle}
Bilirubin is measured by the diazo reaction (van den Bergh reaction). In the classic method, diazotized sulfanilic acid reacts with bilirubin to form colored azobilirubin (azodipyrrole) compounds measured spectrophotometrically at 540--560 nm. Direct (conjugated) bilirubin reacts rapidly in aqueous solution (``direct'' reaction) because it is water-soluble. Unconjugated bilirubin requires a solubilizing agent (caffeine, methanol, or dimethyl sulfoxide) to accelerate the reaction, hence the term ``indirect'' for unconjugated bilirubin. Modern automated analyzers use modifications of the Jendrassik-Gr\'{o}f method (1938), which employs caffeine-benzoate as the accelerator and ascorbic acid to stop the reaction. Some analyzers use direct spectrophotometric methods measuring absorbance at multiple wavelengths (bilirubin absorbs at 460 nm), or dry chemistry methods (Vitros) using a multilayered slide with a mordant that produces a colored complex. In neonates, transcutaneous bilirubinometry provides non-invasive screening but requires serum confirmation for elevated values. Direct spectrophotometric methods avoid the need for diazo reagents. Results are reported in mg/dL or $\mu$mol/L. This test is NOT performed by hematology analyzers.

\section*{History}
Bilirubin was first crystallized by Rudolf Virchow in 1847 and named from the Latin ``bilis'' (bile) and ``ruber'' (red). The diazo reaction was developed by Hijmans van den Bergh in 1913 (the ``van den Bergh reaction'') \cite{vandenbergh1916uber}, who also discovered that some bilirubin reacts ``directly'' while other forms require an accelerator, establishing the distinction between conjugated and unconjugated bilirubin. Gilbert syndrome was described by Augustin Nicolas Gilbert in 1901. Crigler-Najjar syndrome was described in 1952. Dubin-Johnson syndrome was described in 1954. The Jendrassik-Gr\'{o}f method, still the basis of modern bilirubin assays, was published in 1938. Phototherapy for neonatal jaundice was discovered serendipitously at Rochford General Hospital, Essex, UK in 1956 by Sister Jean Ward and formally described by Cremer and colleagues in 1958 \cite{cremer1958influence}.

\section*{Why This Test Is Ordered}
\begin{itemize}
  \item Evaluation of jaundice (yellow discoloration of skin, sclera, and mucous membranes): jaundice becomes visible when total bilirubin exceeds 2.5--3.0 mg/dL
  \item Classification of hyperbilirubinemia as unconjugated (pre-hepatic) vs. conjugated (hepatic/post-hepatic) based on the fractionation pattern, guiding differential diagnosis
  \item Assessment of liver synthetic and excretory function as part of the comprehensive metabolic panel: rising bilirubin with falling albumin and rising INR indicates deteriorating liver function
  \item Screening and monitoring of neonatal jaundice: physiological jaundice appears at 2--5 days of life; pathologic jaundice within first 24 hours requires urgent workup
  \item Diagnosis of hemolytic anemia: elevated unconjugated bilirubin with normal liver enzymes and elevated reticulocyte count, LDH, and low haptoglobin
  \item Evaluation of suspected biliary obstruction: elevated direct bilirubin with elevated ALP and GGT
  \item Monitoring of inherited bilirubin metabolism disorders: Gilbert syndrome, Crigler-Najjar syndrome, Dubin-Johnson syndrome, Rotor syndrome
  \item Prognosis assessment in acute liver failure and cirrhosis: bilirubin is a key component of the MELD score (Model for End-Stage Liver Disease), used for liver transplant prioritization
  \item Monitoring of drug-induced liver injury (DILI) severity: bilirubin >2$\times$ ULN with ALT >3$\times$ ULN defines ``Hy's Law'' cases, which have approximately 10--50\% mortality from acute liver failure
\end{itemize}

\section*{Primary vs Secondary Test}
Total bilirubin is a primary screening test as part of the CMP. Fractionation into direct and indirect bilirubin is a secondary step performed when total bilirubin is elevated, to classify the etiology.

\section*{How This Test Is Performed}
A healthcare professional draws blood from a vein into a serum separator tube (gold-top) or lithium heparin tube (green-top). Important handling considerations: (1) Bilirubin is light-sensitive---the sample must be protected from direct light (wrapped in foil or amber tube) because ultraviolet light photodegrades bilirubin, causing falsely low results. Photodegradation can decrease bilirubin by up to 30--50\% per hour of light exposure. (2) Hemolysis interferes with the diazo reaction and spectrophotometric measurement and should be avoided. (3) The sample is centrifuged and analyzed on an automated chemistry analyzer. Results are available within 30--60 minutes. Neonatal samples require special small-volume handling.

\section*{What Preparation Is Needed}
Fasting for 4--8 hours is recommended. Fasting increases unconjugated bilirubin in patients with Gilbert syndrome (UGT1A1 reduced activity) because fasting reduces hepatic UDP-glucuronic acid availability. In Gilbert syndrome, a 48-hour fast can increase bilirubin to 2--4 mg/dL, and this test was historically used as a diagnostic challenge (now rarely performed). Avoid hemolysis (fist clenching, traumatic venipuncture). Inform the provider of medications that cause hyperbilirubinemia: rifampin (competes for hepatic uptake), probenecid (competes for conjugation), cyclosporine, indinavir and atazanavir (UGT1A1 inhibition---atazanavir causes unconjugated hyperbilirubinemia in 30--50\% of patients, mimicking Gilbert syndrome), and drugs causing cholestasis (estrogens, anabolic steroids). No other special preparation is needed.

\section*{Contraindications and Precautions}
\begin{itemize}
  \item No absolute contraindications. Critical pre-analytical considerations:
  \item (1) Light exposure: bilirubin is extremely photolabile. The sample must be protected from direct light (sunlight, fluorescent/UV light) from the moment of collection through analysis. Wrapping the tube in aluminum foil or using an amber tube is standard practice. Light exposure of 1 hour can decrease bilirubin by 30--50\%.
  \item (2) Hemolysis: hemoglobin from lysed red blood cells interferes with the diazo reaction (falsely low bilirubin) and also contributes its own color, causing spectrophotometric interference.
  \item (3) Lipemia: severe hypertriglyceridemia causes turbidity that interferes with spectrophotometric measurement. The laboratory can ultracentrifuge lipemic samples before analysis.
  \item (4) Drugs directly interfering with bilirubin measurement: salicylates, sulfonamides, and free fatty acids can falsely elevate bilirubin in some diazo methods.
  \item (5) Physiological changes: newborns have higher bilirubin due to immature UGT1A1, shorter red blood cell lifespan, and high hemoglobin turnover. Bilirubin peaks at 3--5 days of life (physiologic jaundice) at 5--12 mg/dL in term infants.
  \item (6) Prolonged standing of uncapped samples may lead to protein precipitation affecting measurement.
\end{itemize}
\section*{Risks}
Minimal risk from venipuncture. Mild pain or bruising resolves in 1--2 days. For neonates, heel-stick capillary sampling is standard and carries minimal risk.

\section*{What Happens After the Test}
Results available within 30--60 minutes. The fractionation pattern determines the direction of investigation: (1) Predominantly unconjugated hyperbilirubinemia (>80\% indirect, direct <15\% of total or <0.3 mg/dL): indicates increased bilirubin production (hemolysis, ineffective erythropoiesis) or impaired hepatic uptake/conjugation (Gilbert, Crigler-Najjar, drug-induced). (2) Predominantly conjugated hyperbilirubinemia (>50\% direct or direct >2 mg/dL): indicates hepatocyte excretory dysfunction (Dubin-Johnson, Rotor, drug-induced, viral/alcoholic hepatitis) or biliary obstruction (choledocholithiasis, stricture, malignancy). (3) Mixed hyperbilirubinemia: diffuse hepatic injury (cirrhosis, hepatitis, sepsis). Critical values: total bilirubin >15 mg/dL in adults warrants urgent evaluation. In neonates, exchange transfusion thresholds depend on gestational age, weight, and risk factors, using nomograms (Bhutani nomogram): phototherapy typically started at 15--18 mg/dL in term infants without risk factors. Kernicterus risk at bilirubin >25--30 mg/dL in term infants. The MELD score uses bilirubin, creatinine, and INR: MELD = 3.78 $\times$ ln[bilirubin] + 11.2 $\times$ ln[INR] + 9.57 $\times$ ln[creatinine] + 6.43; scores >15 qualify for transplant listing.

\section*{Understanding Results}

\subsection*{Below Normal}
\begin{itemize}
  \item Low bilirubin (<0.3 mg/dL) is normal and not clinically significant. Bilirubin has antioxidant properties, and low-normal levels are not associated with disease.
  \item Epidemiological studies suggest mild bilirubin elevation (Gilbert syndrome range) may be protective against cardiovascular disease due to bilirubin's antioxidant effects, but low bilirubin is not a disease marker.
  \item Falsely low results from sample light exposure (photodegradation)---most common cause of unexpectedly low bilirubin in a jaundiced patient.
  \item Hemolysis can cause falsely low bilirubin in diazo methods.
  \item Certain medications (salicylates, barbiturates) can induce UGT1A1 and lower bilirubin.
\end{itemize}

\subsection*{Above Normal}
\begin{itemize}
  \item \textbf{Unconjugated (indirect) hyperbilirubinemia (direct <15--20\% of total):}
  \item \textbf{Hemolysis (increased production):} Hemolytic anemia (autoimmune, hereditary spherocytosis, G6PD deficiency, sickle cell disease, thalassemia, microangiopathic anemia)---elevated reticulocyte count, elevated LDH, low haptoglobin, and unconjugated bilirubin typically 2--5 mg/dL (rarely exceeds 4--5 mg/dL because the liver can conjugate up to 4$\times$ normal bilirubin load). Ineffective erythropoiesis (thalassemia, megaloblastic anemia, myelodysplasia, lead poisoning). Resorption of large hematomas (bruising, retroperitoneal hemorrhage). Massive blood transfusion (>10 units).
  \item \textbf{Impaired hepatic uptake/conjugation:} Gilbert syndrome: most common inherited disorder of bilirubin metabolism, affecting 5--10\% of the population; autosomal recessive, mutation in UGT1A1 promoter (TA repeat) \cite{bosma1995genetic}, causing reduced enzyme activity to approximately 30\% of normal. Mild, fluctuating unconjugated hyperbilirubinemia (1.5--4 mg/dL), normal liver enzymes, normal LDH and haptoglobin, no hemolysis. Exacerbated by fasting, illness, stress, alcohol, and exercise. Benign---no treatment needed. Crigler-Najjar syndrome Type I: complete absence of UGT1A1 activity (autosomal recessive); severe unconjugated hyperbilirubinemia (20--50 mg/dL) from birth; kernicterus risk; requires lifelong phototherapy and ultimately liver transplantation. Crigler-Najjar Type II: severely reduced UGT1A1 activity (autosomal recessive or dominant); bilirubin 6--20 mg/dL; responds to phenobarbital (induces residual enzyme activity). Drug-induced: atazanavir, indinavir, rifampin, probenecid. Hypothyroidism. Neonatal physiologic jaundice (immature UGT1A1). Breast milk jaundice (inhibitor in breast milk, onset after day 7, benign).
  \item \textbf{Conjugated (direct) hyperbilirubinemia (direct >50\% of total or >2 mg/dL):}
  \item \textbf{Biliary obstruction (cholestatic):} Choledocholithiasis (gallstone in common bile duct), pancreatic head carcinoma, cholangiocarcinoma, ampullary carcinoma, PSC, PBC, benign biliary stricture, biliary atresia (neonates). Labs: elevated direct bilirubin, ALP >> ALT, GGT elevated. Imaging (ultrasound, MRCP, ERCP) shows dilated bile ducts.
  \item \textbf{Hepatocellular injury:} Viral hepatitis (HAV, HBV, HCV, EBV, CMV), alcoholic hepatitis, NAFLD/NASH---direct bilirubin proportionally elevated, with AST/ALT elevation often predominant over bilirubin. In severe acute hepatitis, direct bilirubin rises as disease worsens.
  \item \textbf{Inherited conjugated hyperbilirubinemia:} Dubin-Johnson syndrome: rare autosomal recessive; mutation in ABCC2 (MRP2) causing impaired canalicular excretion of conjugated bilirubin. Mild conjugated hyperbilirubinemia (2--5 mg/dL), normal liver enzymes, intensely black liver on gross examination (from epinephrine metabolite accumulation in lysosomes), delayed biliary excretion on hepatobiliary scintigraphy (HIDA scan) but normal visualization of liver. Benign---no treatment. Rotor syndrome: rare autosomal recessive; mutation in SLCO1B1 and SLCO1B3 (OATP1B1/1B3) causing impaired hepatic uptake and storage of bilirubin. Similar presentation to Dubin-Johnson but liver is not pigmented and HIDA scan shows no liver visualization. Also benign.
  \item \textbf{Sepsis and TPN:} Sepsis causes cholestasis (``ICU jaundice'') from cytokine effects on bile transporters. TPN-induced cholestasis from prolonged parenteral nutrition (especially in infants and those on long-term TPN). Drug-induced cholestasis from estrogens, anabolic steroids, erythromycin, amoxicillin-clavulanate, chlorpromazine, cyclosporine.
  \item \textbf{Critical bilirubin:} >15 mg/dL in adults; >18--20 mg/dL in neonates (phototherapy threshold).
\end{itemize}

\section*{Normal Results}
\begin{itemize}
  \item \textbf{Total bilirubin (adults):} 0.3--1.2 mg/dL (5--21 $\mu$mol/L)
  \item \textbf{Direct (conjugated) bilirubin (adults):} 0.0--0.3 mg/dL (0--5 $\mu$mol/L)
  \item \textbf{Indirect (unconjugated) bilirubin (adults):} 0.2--0.9 mg/dL (calculated as total -- direct)
  \item \textbf{Newborns (term, day 1):} Total <6 mg/dL; Day 3--5: <12--15 mg/dL
  \item \textbf{Newborns (preterm, day 1):} Lower phototherapy thresholds; weight and gestational age-based nomogram
  \item \textbf{Exchange transfusion threshold (term, healthy):} 20--25 mg/dL (depending on age in hours and risk factors)
  \item \textbf{Jaundice visible:} Total bilirubin >2.5--3.0 mg/dL
  \item \textbf{Critical value (adult):} Total bilirubin >15 mg/dL
  \item \textbf{MELD score calculation:} Uses bilirubin, creatinine, INR, and sodium
  \item \textbf{Conversion:} 1 mg/dL = 17.1 $\mu$mol/L
\end{itemize}

\section*{Follow-up Tests}
\begin{itemize}
  \item \textbf{Direct (conjugated) bilirubin:} Essential to fractionate when total is elevated. Direct >2 mg/dL or >50\% of total indicates conjugated hyperbilirubinemia.
  \item \textbf{CBC with reticulocyte count:} Evaluate for hemolysis as cause of unconjugated hyperbilirubinemia.
  \item \textbf{LDH and haptoglobin:} LDH elevated and haptoglobin low in hemolysis.
  \item \textbf{Peripheral blood smear:} Look for spherocytes, schistocytes, sickle cells.
  \item \textbf{Liver function panel (ALT, AST, ALP, GGT, albumin, INR):} Complete assessment of liver injury pattern and synthetic function.
  \item \textbf{Abdominal ultrasound with Doppler:} First-line imaging for biliary dilation, gallstones, liver texture, portal hypertension.
  \item \textbf{MRCP or ERCP:} Detailed biliary imaging for suspected obstruction.
  \item \textbf{Viral hepatitis serologies:} HAV IgM, HBsAg, anti-HBc IgM, anti-HCV.
  \item \textbf{Anti-mitochondrial antibody (AMA):} For PBC.
  \item \textbf{Direct Coombs test:} If autoimmune hemolytic anemia suspected.
  \item \textbf{Genetic testing for UGT1A1:} For Gilbert or Crigler-Najjar syndrome.
\end{itemize}

\section*{References}
\bibliographystyle{unsrtnat}
\bibliography{references}

\clearpage
\section*{Regulatory Status and Authorization Date}
This chapter describes a test category, not a specific commercial product. FDA, CDSCO, EU IVDR, and MHRA approval or registration dates therefore cannot be assigned without the assay manufacturer, product name, instrument, and intended use. For laboratory-developed tests, an individual agency approval date may not exist. See the Preface for the interpretation of regulatory status.

\clearpage
